\section{Kategori dan Tujuan Serangan Kriptografi}

Kriptanalisis tidak hanya terbatas pada upaya menebak kunci, tetapi mencakup berbagai strategi untuk mengganggu keamanan informasi. Serangan dikategorikan berdasarkan cara penyerang berinteraksi dengan sistem dan apa yang ingin dicapai pada akhirnya.

\subsection{Serangan Pasif vs Serangan Aktif}

Interaksi penyerang dengan saluran komunikasi dapat dibedakan menjadi dua kategori utama:

\begin{enumerate}
    \item \textbf{Serangan Pasif (\textit{Passive Attack})} \\
    Penyerang hanya berperan sebagai pengamat. Ia memantau, merekam, dan menganalisis transmisi data tanpa mengubah satu bit pun dari pesan yang dikirim.
    \begin{itemize}
        \item \textbf{Eavesdropping (Penyadapan)}: Mengintercept paket data di jaringan atau mendengarkan percakapan. Tujuan utamanya adalah melanggar \textit{confidentiality}.
        \item \textbf{Traffic Analysis (Analisis Lalu Lintas)}: Meskipun isi pesan terenkripsi, penyerang menganalisis pola komunikasi—siapa yang mengirim, kapan, berapa sering, dan berapa besar datanya. Informasi ini seringkali cukup untuk mengungkap struktur organisasi atau rencana operasi.
    \end{itemize}
    Karakteristik utama serangan pasif adalah \textbf{sulit dideteksi}, karena tidak ada perubahan pada data yang mengalir di jaringan. Strategi terbaik untuk menghadapinya adalah melalui enkripsi yang kuat.

    \item \textbf{Serangan Aktif (\textit{Active Attack})} \\
    Penyerang secara sengaja melakukan modifikasi pada aliran data atau menyuntikkan data palsu ke dalam sistem. Serangan ini melanggar \textit{integrity} dan \textit{availability}.
    \begin{itemize}
        \item \textbf{Masquerading (Pemalsuan Identitas)}: Penyerang berpura-pura menjadi pengguna sah untuk mendapatkan akses atau mengirim pesan palsu.
        \item \textbf{Modification (Modifikasi)}: Mengubah isi pesan yang dicegat (misal: mengubah nominal transfer bank dari Rp 1.000.000 menjadi Rp 10.000.000) sebelum meneruskannya ke penerima.
        \item \textbf{Replay Attack (Serangan Pengulangan)}: Menangkap pesan valid (misal: permintaan login) dan mengirimkannya kembali di waktu lain untuk menipu sistem.
        \item \textbf{Denial of Service (DoS)}: Membanjiri sistem dengan permintaan palsu sehingga pengguna sah tidak dapat mengakses layanan.
    \end{itemize}
    Serangan aktif lebih mudah dideteksi karena menyebabkan anomali pada data atau gangguan layanan, namun dampak kerusakannya jauh lebih destruktif.
\end{enumerate}

\subsection{Hierarki Tujuan Kriptanalisis}

Tergantung pada motifnya, seorang kriptanalis menetapkan target keberhasilan yang berbeda-beda:

\begin{itemize}
    \item \textbf{Pemulihan Plainteks (\textit{Plaintext Recovery})} \\
    Tujuan jangka pendek: Membaca isi pesan tertentu. Penyerang hanya ingin tahu apa yang dikatakan dalam satu pesan spesifik tanpa harus tahu kuncinya.

    \item \textbf{Pemulihan Kunci (\textit{Key Recovery})} \\
    Tujuan jangka panjang: Menemukan kunci rahasia. Ini adalah serangan paling fatal, karena memberikan akses ke seluruh komunikasi masa lalu dan masa depan yang menggunakan kunci tersebut.

    \item \textbf{Pemalsuan Pesan (\textit{Forgery})} \\
    Tujuan manipulatif: Menciptakan cipherteks yang tampak sah. Penyerang tidak perlu tahu kunci, tetapi ia ingin sistem penerima menerima pesan palsu sebagai pesan asli yang valid.

    \item \textbf{Pembedaan (\textit{Distinguishing Attack})} \\
    Tujuan teoretis: Membuktikan bahwa sebuah cipher memiliki kelemahan. Penyerang mencoba membuktikan bahwa cipherteks tidak sepenuhnya acak, sehingga ia dapat membedakan output cipher dari deretan bit acak murni.
\end{itemize}

\subsection{Korelasi dengan Pilar CIA}

Setiap kategori serangan secara langsung mengancam salah satu pilar keamanan informasi:
\begin{itemize}
    \item \textbf{Confidentiality} $\rightarrow$ Terancam oleh serangan pasif dan pemulihan kunci.
    \item \textbf{Integrity} $\rightarrow$ Terancam oleh serangan aktif berupa modifikasi dan pemalsuan.
    \item \textbf{Availability} $\rightarrow$ Terancam oleh serangan aktif berupa Denial of Service (DoS).
\end{itemize}
